\section{Conclusion and Future Work}
\label{sec:conclusion}

In this paper, guided by an empirically-driven research paradigm, we exposed the \textbf{Overfitting-Optimizer Paradox} in adaptive model merging. We demonstrated that unconstrained gradient-based test-time adaptation (such as AdaMerging under Adam GD) is highly susceptible to transductive overfitting on unlabeled target streams. This causes the optimizer to construct jagged, physically unrealistic merging coefficient profiles that fail to generalize on held-out distributions. 

To resolve this issue, we introduced \textbf{PolyMerge}, which projects the merging coefficient search space onto a continuous, low-dimensional polynomial subspace of normalized layer depth. By optimizing only $d+1$ parameters per task, PolyMerge hard-constrains the search space to mathematically enforce physical depth-wise smoothness and eliminate high-frequency optimization noise. Our exhaustive optimization sweeps within our calibrated Vision Transformer simulation environment across multiple benchmarks, optimizers, complexity degrees, and random seeds confirmed that PolyMerge completely prevents transductive overfitting. Notably, PolyMerge ($d=2$, Adam) achieved a state-of-the-art average accuracy of \textbf{86.57\%}, dramatically outperforming both uniform Task Arithmetic and unconstrained AdaMerging while mapping a classic bias-variance curve peaking at the quadratic degree.

Several exciting avenues of future research remain. First, we plan to scale PolyMerge to massive Large Language Models (LLMs) with dozens or hundreds of layers, where unconstrained adaptation is expected to suffer from even more severe dimensionality-driven overfitting. Second, we aim to explore other continuous, low-frequency representation spaces, such as low-frequency Discrete Cosine Transforms (DCT) or piecewise cubic splines, to model more complex but still globally regularized layer transitions. Finally, we hope to conduct a formal geometric analysis of the adaptive merging landscape to study the relationship between loss landscape flatness, parameter constraints, and out-of-distribution generalization.
